\section{A simultaneous cover for two actual quadratic-domain powers}
\label{bk-section}

Work in the fixed degree-four field
$K=\mathbb Q(i,\sqrt{-3})$. Write
\[
 \zeta=(1+\sqrt{-3})/2,\qquad
 \alpha=(\sqrt3-i)/2,\qquad
 \overline\alpha=(\sqrt3+i)/2,
 \qquad \alpha^2=\overline\zeta.
\]
Complex conjugation acts on coefficients and leaves the variable $t$
fixed. Put
\begin{gather}
 a_0=t^2-1,\quad b_0=2t+1,\quad c_0=t^2+2t,\quad
 U_0=t^2+t+1,\quad A_0=a_0b_0,\label{bk-parameter}\\
 E=A_0+\zeta U_0^2,\qquad G=A_0+iU_0c_0.
 \label{bk-quadratic-elements}
\end{gather}
Then $c_0^2-A_0=U_0^2$ and
$E\overline E=G\overline G=F(a_0,b_0)$.

\begin{lemma}[The full branch inventory]\label{bk-branches}
Set
\[
 B_1=c_0-\alpha U_0,\quad B_2=c_0+\alpha U_0,\quad
 B_3=c_0-\overline\alpha U_0,\quad
 B_4=c_0+\overline\alpha U_0.
\]
These four quadratics have eight distinct simple roots, none at infinity,
and
\begin{equation}\label{bk-four-factorizations}
 E=B_1B_2,\quad\overline E=B_3B_4,\quad
 G=B_1B_4,\quad\overline G=B_2B_3.
\end{equation}
\end{lemma}
\begin{proof}
The factorizations use $A_0=c_0^2-U_0^2$,
$1-\zeta=\overline\zeta$, $\overline\alpha-\alpha=i$ and
$\alpha\overline\alpha=1$. Each quadratic is $c_0-sU_0$, with
$s\in\{\alpha,-\alpha,\overline\alpha,-\overline\alpha\}$.
These four constants are distinct and different from one. Its leading
coefficient is $1-s$, and its discriminant is
\[
 (2-s)^2+4s(1-s)=4-3s^2\ne0.
\]
Here $s^2$ is $\zeta$ or $\overline\zeta$, and the displayed
discriminant has norm $13$ from $\mathbb Q(\sqrt{-3})$.
A root shared by two quadratics would force $c_0=U_0=0$.
But $U_0-c_0=1-t$ and $U_0(1)=3$. There is no common finite root,
and the nonzero leading coefficients exclude infinity.
\end{proof}

\begin{theorem}[Exact component count and genus]\label{bk-degree-genus}
For an integer $g\ge2$ and constants $\eta_E,\eta_G\in K^*$, consider
\begin{equation}\label{bk-cover}
 Y^g=\eta_E\frac E{\overline E},\qquad
 Z^g=\eta_G\frac G{\overline G}.
\end{equation}
Take the smooth projective normalizations of these equations.
Put $\delta=\gcd(g,2)$. There are exactly $\delta$
geometric components. Each has degree $g^2/\delta$ over
$\mathbb P^1_t$ and genus
\begin{equation}\label{bk-genus}
 1+\frac{3g^2-4g}{\delta}.
\end{equation}
Thus the odd-exponent curve is geometrically connected. At an even
exponent there are two geometric components; their degrees are
$g^2/2$, not $g^2$.
\end{theorem}
\begin{proof}
Over an algebraic closure, the two functions in
\eqref{bk-cover} have, on the two roots of each successive $B_j$,
the valuation vectors
\begin{equation}\label{bk-valuation-matrix}
 (1,1),\quad(1,-1),\quad(-1,-1),\quad(-1,1).
\end{equation}
A relation with exponents $u,v$ modulo $g$ that is a $g$th power must
satisfy $u+v\equiv u-v\equiv0\pmod g$. This gives one relation
when $g$ is odd and two when $g$ is even. They are sufficient: the
nontrivial even relation is $(u,v)=(g/2,g/2)$ and uses
\[
 \frac E{\overline E}\frac G{\overline G}
   =\left(\frac{B_1}{B_3}\right)^2.
\]
Constants have $g$th roots over the algebraic closure. Kummer theory
therefore gives subgroup order, and function-field degree,
$g^2/\delta$. The two equations have total generic degree $g^2$ and
split into exactly $\delta$ geometric components of that degree.

At each of the eight points of Lemma~\ref{bk-branches}, both functions
have valuation $1$ or $-1$. After taking $g$th roots of local units,
the local extension is generated by one $g$th root of a uniformizer.
Its inertia order on each component is exactly $g$. At infinity both
functions have order zero and are local units, and there is no
ramification there or at any other point. Set $d=g^2/\delta$.
Characteristic-zero Riemann--Hurwitz gives
\[
 2\mathfrak g-2=-2d+8d(1-1/g),
\]
which proves~\eqref{bk-genus}.
\end{proof}
For the pure unit constants below, the two even components are already
defined over $K$. Indeed $\eta_E\eta_G\in\mu_6$, and its square
roots lie in $\mu_{12}\subset K$. The components are distinguished by
\begin{equation}\label{bk-even-components}
 \frac{(YZ)^{g/2}}{B_1/B_3}
   =\sqrt{\eta_E\eta_G}\quad\hbox{or}\quad
    -\sqrt{\eta_E\eta_G}.
\end{equation}
This arithmetic component assertion is restricted to those constants;
the preceding theorem asserts the geometric count for arbitrary
$\eta_E,\eta_G\in K^*$.

\begin{theorem}[Actual seeds lift to finitely many unit twists]
\label{bk-actual-lift}
Let $a,b$ be positive coprime integers with $M=U^2$ for a positive
integer $U$ and $F=Q^g$, where $Q>1$ and $g\ge2$ are integers.
Then this seed gives a nonbranch $K$-rational point on one of the at
most six curves~\eqref{bk-cover} with
$\eta_E\in\mu_3$, $\eta_G\in\mu_2$, at base coordinate
\begin{equation}\label{bk-actual-parameter}
 t=(a+U)/b>1.
\end{equation}
For odd $g$ only the three $\eta_E$ twists are required. If
$\gcd(g,6)=1$, both constants can be absorbed: every such seed lifts
to the single curve with $\eta_E=\eta_G=1$, of degree $g^2$ and
genus $1+3g^2-4g$.
\end{theorem}
\begin{proof}
The equality $U^2-(a-b)^2=3ab>0$ gives $t>1$.
For $x=a/b$ one has $t=x+\sqrt{x^2+x+1}$ and
$t^2-2xt-x-1=0$. Thus $a_0/b_0=x$ and $b_0=2t+1>0$.
With $k=b_0/b>0$,
\begin{equation}\label{bk-positive-scaling}
 (a_0,b_0,c_0)=k(a,b,c),\qquad U_0=kU,
 \qquad E=k^2(ab+U^2\zeta),\quad G=k^2(ab+iUc).
\end{equation}
The sign of $U_0=kU$ is fixed by $U_0=t^2+t+1>0$.

The two actual primitive UFD decompositions give
\[
 ab+U^2\zeta=u_Ew_E^g,\qquad ab+iUc=u_Gw_G^g,
 \qquad\operatorname N(w_E)=\operatorname N(w_G)=Q.
\]
Consequently $Y=w_E/\overline w_E$ and $Z=w_G/\overline w_G$
satisfy~\eqref{bk-cover} with $\eta_E=u_E^{-2}\in\mu_3$ and
$\eta_G=u_G^{-2}\in\mu_2$. The positive rational factor $k^2$
cancels from both ratios. None of $E,\overline E,G,\overline G$
vanishes at this point, since their norms are $k^4F>0$. Hence the
point is smooth and nonbranch and lifts to the normalization.

For odd $g$, raising to $g$ permutes $\mu_2$, so $\eta_G$ can be
absorbed into $Z$. If $3\nmid g$, it also permutes $\mu_3$, so
$\eta_E$ can be absorbed into $Y$. These changes take place in the
fixed field $K$ and require no varying cyclotomic field.
\end{proof}

\paragraph{Verification and remaining scope.}
The ordinary proof computes the exact relation module, including the
even-exponent defect, before asserting a degree. Its geometric inputs
are the same characteristic-zero Kummer and Riemann--Hurwitz facts
used in the preceding reviewed cover constructions. The supplemental
exact replay checks the four polynomial factorizations, their norm
identity, four discriminants and six nonzero pairwise resultants in
$\mathbb Q(\alpha)$, and the relation kernels at $2\le g\le120$.
It also checks $52$ actual square-first seeds with
$1\le a\le b\le300$, their complete trial factorizations, and $74$
integer inverse instances with $c\le600$. These finite checks do not
establish rational-point completeness.

This is a further cover to which actual points lift, not a birational
replacement of the earlier double-oriented curve and not an integer
inverse for arbitrary $K$-points. The larger genus may be useful for
fixed-exponent rational-point calculations; it supplies no calculation
or height bound uniform in $g$. The even-exponent geometric components
exist despite the already proved emptiness of the positive seed locus.
The odd second exponent and the ramified first branch remain open.
No complete Lean formalization of these geometric assertions is claimed.
